% !TeX encoding = UTF-8
% !TeX program = xelatex
% !TeX spellcheck = fr_CA
%%%%%%%%%%%%%%%%%%%%%%%%%%%%%%%%%%%%%%%%%%%%%%%%%%%%%%%%%%%%%%%%%%%%%%%%%%%%%%%%
%  Exemple minimal BeamerJS : le point de départ d'une nouvelle présentation.
%  Exemple d'utilisation de la classe BeamerJS de Jérôme Soucy.
%  Contrairement à BeamerJS.cls, ce fichier peut être copié et modifié
%  librement, sans condition : c'est le point de départ de vos présentations.
%
%  Pour tout le reste, voir demonstration.tex.
%%%%%%%%%%%%%%%%%%%%%%%%%%%%%%%%%%%%%%%%%%%%%%%%%%%%%%%%%%%%%%%%%%%%%%%%%%%%%%%%
\documentclass{BeamerJS}

\title{Titre de l'exposé}
\subtitle{Sous-titre facultatif}
\author{Prénom Nom}
\date{\today}

\begin{document}

\begin{frame}[plain,noframenumbering]
	\titlepage
\end{frame}

\section{Introduction}

\begin{frame}{Un premier résultat}
	\begin{thm}[Pythagore]
		Dans un triangle rectangle d'hypoténuse $c$, on a $a^2+b^2=c^2$.
	\end{thm}

	\begin{apoint}
		Remplacez ce contenu par le vôtre.
	\end{apoint}
\end{frame}

\end{document}
